
LLM-guided theorem provers achieve substantially higher success rates on miniF2F olympiad problems than pure automated provers, yet no prior work quantifies which mechanisms drive this advantage. DeepSeek-Prover-V2 reports 88.9\% success on miniF2F-test (2025), while automated provers achieve much lower success rates. Thor \citep{Thor_2022} demonstrated that hybrid LLM+hammer systems produce 8.2\% unique solutions where neither component succeeds alone, revealing complementary strengths. Without mechanistic understanding, we cannot predict when synergy occurs or design targeted improvements.

Understanding LLM advantage mechanisms is critical for designing hybrid systems that combine LLM semantic understanding with automated prover formal reliability. If the gap originates from natural language hint parsing, we should invest in NL preprocessing and semantic tactic libraries. If it stems from long-range proof search, we should prioritize context window scaling. Existing work reports that LLMs outperform automated provers and how much they outperform, but not which mechanisms drive the advantage.

We hypothesize the gap decomposes into three testable mechanisms: (1) natural language understanding (predicted 60\% contribution)---LLMs exploit informal hints in problem statements that automated provers ignore, (2) proof depth capability (30\%)---LLMs maintain context across multi-step proofs where automated provers timeout at shallow depths, and (3) corpus pattern matching (10\%)---LLMs learn human proof tactic distributions from Mathlib training data. Each mechanism is falsifiable via controlled ablation: remove NL hints, filter to shallow proofs, compare against random Mathlib sampling.

Our key finding: natural language understanding is the dominant mechanism, contributing ~60\% of LLM advantage (29.5 percentage point drop when NL hints removed, p<$10^{-9})$, while proof depth contributes less than 8\% (below our 5\% falsification threshold). The depth result derives from infrastructure validation using simplified problems (96.3\% solvable with $\leq 3$ tactics, unrealistic for olympiad mathematics where literature reports median=9 tactics) and requires revalidation on real miniF2F data. This finding shifts design priorities from architectural scaling to NL-formal integration. For hybrid systems, mechanistic attribution suggests routing NL-rich problems to LLMs and formal-only goals to automated provers.

\textbf{Contributions.} Building on this mechanistic hypothesis, we:

\begin{enumerate}
\item Establish the first measured baseline for pure automated provers (lean-auto) on miniF2F Lean 4 subset: 15.6\% [11.5\%, 20.3\%] (N=244), validating our 15\% prediction and enabling controlled LLM comparisons.
\end{enumerate}

\begin{enumerate}
\item Quantify NL understanding contribution via ablation study: removing docstring/comment hints drops LLM success from 62.3\% to 32.8\% ($\Delta =29.51$\% [20.90\%, 38.11\%], p<$10^{-9})$, directly validating ~60\% attribution.
\end{enumerate}

\begin{enumerate}
\item Reject the proof depth hypothesis provisionally: post-hoc filtering to shallow proofs ($\leq 3$ tactics) drops success by only 3.7\% [1.6\%, 6.1\%] (below our 5\% gate threshold). This result derives from infrastructure validation using simplified problems and requires revalidation on real olympiad-difficulty data.
\end{enumerate}

\begin{enumerate}
\item Validate tactic budget control framework: tactic count coefficient of variation CV=0.36 (36\% << 100\% threshold) enables fair LLM vs automated prover comparisons by controlling computational confounds (budget=15 captures 90.6\% of baseline strategies).
\end{enumerate}

Our mechanistic approach differs from prior systems work (Thor, LeanCopilot, DeepSeek-Prover-V2) in focus: we explain LLM advantage rather than maximize success rates.
